\section{NFT Liquidity Protocol Post-Quantum}
\label{sec:nft-liquidity-pq}

The NFT Liquidity Protocol --- the term \Zoo{} coined in the
October 31, 2021 whitepaper \cite{zoo-2021} --- turns \NFT{}s into
productive collateral inside a yield-bearing system. In 3.0 the
loan, repayment, and liquidation primitives are all post-quantum
authenticated.

\subsection{Why this matters specifically for NFT collateral}

An NFT loan is a long-lived bilateral commitment: the borrower locks
the \NFT{}, the lender extends a loan, and the agreement persists
for the loan duration (months to years). A classical signature on
that agreement is exactly the kind of long-lived hot-path credential
that harvest-now-decrypt-later most threatens. If the agreement
expires post-Q-day, the holder of the original loan agreement may
discover their authority has been retroactively forged. 3.0 closes
that gap.

\subsection{Loan agreement schema (3.0)}

\begin{lstlisting}[language=C++,caption={3.0 NFT loan agreement schema}]
struct NftLoanAgreement3 {
    Hash      nft_id;                  // ZRC-721 token reference
    Address   borrower;
    Address   lender;
    Amount    principal;
    Rate      apr;
    Duration  term;
    Hash      collateral_appraisal;    // F-Chain FHE handle (optional)
    MlDsa::Sig borrower_sig;           // ML-DSA-65 signature
    Corona::Sig lender_sig;          // Corona signature, threshold capable
    Hash      l1_anchor;               // Quasar 3.0 cert reference
};
\end{lstlisting}

The borrower signs with $\MLDSA{}\text{-}65$ to bind individual
identity. The lender signs with \Corona{}, threshold-capable so
that pooled-lender markets (lender DAOs, lending pools) can present
a single threshold signature standing in for many participants. Both
signatures are required for the loan to bind.

\subsection{Confidential appraisal}

The optional \texttt{collateral\_appraisal} field is a handle to an
F-Chain \FHE{} ciphertext holding the appraised value of the \NFT{}.
A confidential market can run liquidation logic over the ciphertext
without exposing the appraised value to public observers. This is
useful for high-value collections where revealing appraisal is itself
information leakage. The 3.0 \FHE{} substrate is the same one used
elsewhere in the stack (LP-013, Confidential ERC-20 LP-067), so
appraisals can be composed with confidential balance transfers
without breaking PQ guarantees.

\subsection{Liquidation flow}

Liquidation is the highest-stakes path in the protocol because it
transfers the \NFT{} away from the borrower. The 3.0 liquidation
flow:

\begin{enumerate}[leftmargin=1.4em]
  \item Liquidation oracle (an A-Chain attested process) reads the
        appraisal-vs-debt threshold.
  \item Oracle issues a liquidation event signed under threshold
        $\MLDSA{}$.
  \item Loan contract verifies the oracle's threshold signature
        against the on-chain oracle set's public key.
  \item Contract transfers \NFT{} to lender (or auctions to the
        lending pool, depending on agreement type).
  \item Liquidation receipt is anchored to \Lux{} \Quasar{} 3.0 cert
        \cite{lp-020}, providing a triple-cert-final record of the
        transfer.
\end{enumerate}

A classical-only oracle compromise (a future Q-day adversary forging
the oracle's classical signature) cannot trigger liquidation in 3.0
because the liquidation contract does not honour classical-only
signatures.

\subsection{Repayment flow}

Repayment closes the loan and unlocks the \NFT{}. The borrower
submits a repayment transaction signed with $\MLDSA{}$; the contract
verifies the signature, transfers the principal plus accrued interest
to the lender, and releases the \NFT{} from escrow. The repayment
receipt is anchored to L1 cert and is independently verifiable
post-Q-day.

\subsection{Pooled lending}

Pooled lender markets aggregate many lenders behind a single
\Corona{} threshold key. The threshold structure means individual
lender shares can be added or removed without re-issuing every
outstanding loan: the pool's threshold public key is stable across
lender membership changes inside the agreed parameter set. When
the pool needs to rotate its lattice parameters (a security-driven
upgrade), it emits a one-time PQ-signed parameter-update event, and
all in-flight loans pin to the parameter epoch in which they were
signed.

\subsection{Migration of pre-3.0 loan agreements}

\begin{enumerate}[leftmargin=1.4em]
  \item All open 2.x loan agreements are walked by an automated
        migration job.
  \item For each open loan, both counterparties are notified to
        co-sign a 3.0 loan agreement that supersedes the 2.x form,
        carrying the same economic terms.
  \item Loans whose counterparties do not co-sign within the
        90-day sunset window are flagged ``classical-only'' on the
        UI; liquidation against these loans requires an additional
        consensus check.
  \item Newly-issued loans after the 3.0 cutover are PQ-only.
\end{enumerate}

This is the same sunset pattern used for governance keys
(\S\ref{sec:dao-migration}) and zLLM commitments
(\S\ref{sec:zllm-pq}).
